% =========================================================================================
% Chapter 2 -- Project Environment and Presentation
% =========================================================================================
\section{Project Environment and Presentation}
\label{sec:environment}

\subsection{Host Organization}
\label{subsec:host}

\rmq{Describe the host organisation: its activity, size, position in its sector, and the
department or team the internship took place in. Half a page to a page. If the project was
carried out within the school rather than a company, describe the laboratory or programme
instead.}

\subsection{Internship Environment}
\label{subsec:internship-env}

The project was carried out as an individual end-of-studies internship over a period of
three months, with the scope of a minimum viable product (\gls{mvp}) rather than a
deployable clinical product. All
design, implementation, experimentation and evaluation described in this report were
performed by the author, with periodic supervision and review.

The technical environment reflects that scope. Development was carried out in Python,
using PyTorch for the imaging models and scikit-learn together with XGBoost for the
structured-data model. Training was performed on Google Colab GPU instances, with local
execution reserved for data preparation, analysis and inference. \rmq{State the supervision
arrangement: frequency of meetings, review format, and any clinical input received.}

The intermittent availability of GPU hardware was a genuine constraint on the experimental
programme rather than an incidental inconvenience. It limited the number of training runs
that could be spent on any single question, and consequently shaped the methodology
described in Section~\ref{subsec:methodology}: candidate changes were, wherever possible,
estimated analytically from measurements already obtained before being committed to a
training run. This constraint is discussed further in Chapter~\ref{sec:conclusion}.

\subsection{Project Context}
\label{subsec:project-context}

The clinical problem was set out in Chapter~\ref{sec:introduction}: point-of-care
ultrasound is fast and available, but operator-dependent, and a single imaging finding
rarely determines a decision on its own. Existing automated approaches, reviewed in
Chapter~\ref{sec:sota}, address one organ and one question at a time.

The project therefore targets the layer above the individual model. Its contribution is
not a new segmentation architecture or a new classifier, but a system in which several
perception models cooperate with structured clinical data under an explicit account of
what each of them can and cannot support. That framing determines what counts as success:
a module is useful to this system when its output is accompanied by a trustworthy
statement of its own reliability, not merely when its headline metric is high.

\subsection{POCUS-Emergency}
\label{subsec:pocus-emergency}

POCUS-Emergency is a modular multi-agent assistant for emergency ultrasound and clinical
decision support. It accepts an ultrasound examination together with whatever structured
clinical data is available for the same patient, and returns a ranked differential
diagnosis, an explicit account of the uncertainty attached to it, and a decision on
whether the case should be escalated for further examination.

The system covers four perception modules --- triage, lung, cardiac and gallbladder ---
and the reasoning layer that combines them. A vascular / deep vein thrombosis (\gls{dvt})
module was considered
during planning and deliberately excluded from the scope of the internship, in order to
complete the four modules above to a defensible standard rather than five to a partial
one. The examination types the Ultrasound Agent does not support return an explicit
\textit{not supported} status rather than a low-confidence guess, for the reasons given in
Section~\ref{subsec:architecture}.

\subsection{System Architecture}
\label{subsec:architecture}

The system is organised as three agents with strictly separated responsibilities,
summarised in Tab.~\ref{tab:agents}. The separation is not incidental: it is the mechanism
by which the fourth difficulty of Section~\ref{subsec:problem} --- unwarranted confidence
--- is controlled.

\begin{table}[H]
  \centering
  \small
  \begin{tabularx}{\linewidth}{lXX}
    \toprule
    \textbf{Agent} & \textbf{Input} & \textbf{Output} \\
    \midrule
    Triage Agent & Chief complaint, vital signs, age, red-flag symptoms &
      Urgency level with calibrated confidence \\
    \addlinespace[2pt]
    Ultrasound Agent & Ultrasound image or clip, routed by organ &
      Findings with calibrated confidence and attention maps \\
    \addlinespace[2pt]
    Clinical Reasoning Agent & Both of the above, plus laboratory results and history &
      Ranked differential diagnosis and escalation decision \\
    \bottomrule
  \end{tabularx}
  \caption{Responsibilities of the three agents.}
  \label{tab:agents}
\end{table}

Each component rests on a distinct technology, and the report names them precisely rather
than describing the system uniformly as ``artificial intelligence''. Tab.~\ref{tab:tech}
gives the correspondence.

\begin{table}[H]
  \centering
  \small
  \begin{tabularx}{\linewidth}{lX}
    \toprule
    \textbf{Component} & \textbf{Technology} \\
    \midrule
    Triage Agent            & Gradient-boosted decision trees (XGBoost
                              \cite{chen2016xgboost}) with post-hoc probability
                              calibration \\
    Lung module             & \gls{cnn} classifier, multi-label, EfficientNet-B0 backbone
                              \cite{tan2019efficientnet} \\
    Cardiac module          & U-Net segmentation \cite{ronneberger2015unet} with an
                              EfficientNet-B0 encoder \\
    Gallbladder module      & \gls{cnn} classifier, EfficientNet-B0 backbone \\
    Explainability          & \gls{gradcam} \cite{selvaraju2017gradcam} \\
    Clinical state          & Deterministic Python; no model \\
    Knowledge grounding     & \gls{rag} over a curated corpus \cite{lewis2020rag} \\
    Differential reasoning  & Medical \gls{llm} \cite{chen2024huatuogpto1} \\
    Escalation              & Rule-based policy; deterministic \\
    \bottomrule
  \end{tabularx}
  \caption{Technology underlying each component of the system.}
  \label{tab:tech}
\end{table}

\subsubsection{Triage Agent}

The Triage Agent reads the patient record rather than the image. From the chief complaint,
vital signs, age and available red-flag indicators it estimates an urgency tier --- low,
medium or high --- together with a confidence. It operates on structured data only and
never sees an ultrasound image.

Two properties matter for the system as a whole. First, its confidence is calibrated, so
that a stated probability can be compared meaningfully against the confidence attached to
an imaging finding. Second, its output is deliberately coarse: three tiers rather than the
five-level scales used in hospital triage systems, because the reasoning layer uses the
tier to decide whether to escalate, not to assign a waiting time. The construction and
evaluation of this model are described in Section~\ref{subsec:triage-agent}.

\subsubsection{Ultrasound Agent}

The Ultrasound Agent is a router with organ-specific models behind it. Given an
examination and the organ examined, it dispatches the input to the corresponding module
and returns findings in a single shared output format, regardless of which module
produced them. Each finding carries a calibrated confidence, an indication of whether that
calibration was actually available, and a \gls{gradcam} map relating the prediction to
image regions the clinician can verify.

\begin{keystatement}
The Ultrasound Agent reports findings and does not assign urgency.
\end{keystatement}

This distinction is deliberate and is enforced throughout the implementation. A
radiologist records ``multiple B-lines, small consolidation, no pleural effusion''; the
inference to pulmonary oedema, and the decision about how urgently to act, require vital
signs, laboratory results and history that the imaging model never sees. Allowing an
imaging model to emit an urgency verdict would place a clinical decision inside a
component that lacks the information to make it, and would duplicate that logic in two
places within the system.

The same reasoning governs unsupported examinations. An organ for which no module was
built returns an explicit \textit{not supported} status in the ordinary output format,
rather than being silently omitted or answered with a low-confidence guess. A missing
module and a negative finding are different statements, and the reasoning layer must be
able to tell them apart.

\subsubsection{Clinical Reasoning Agent}

The Clinical Reasoning Agent occupies the position of the senior physician who receives
the nurse's rapid assessment and the imaging report, weighs both against the laboratory
results and the patient's history, and reaches a conclusion. It is the only component
permitted to decide urgency or recommend action.

\begin{keystatement}
A second separation applies within the reasoning agent itself: deterministic rules decide
what is safe, and the language model explains.
\end{keystatement}

The safety-critical determinations --- which tests are absent, which model limitations
apply, whether the agents disagree, whether escalation is warranted --- are computed by
deterministic rules before the language model is consulted, and the model's output is
checked against the structured state afterwards. The language model reasons over evidence
that has already been constrained; it does not decide whether the case is safe. Asking a
language model directly what should be done would place the least predictable component in
charge of the most consequential decision. This design is described in
Chapter~\ref{sec:reasoning}.

\subsection{Development Methodology}
\label{subsec:methodology}

\subsubsection{CRISP-DM and how it was adapted}

The project followed \gls{crisp}, the reference process model for data mining and data
science work. It was chosen over a software-engineering methodology because the central
risks here are not implementation risks. The difficulty in this project was rarely writing
the code; it was establishing what the data actually represented and whether a measured
result could be trusted. \gls{crisp} places those questions in named phases with explicit
feedback loops, which matches how the work behaved.

Tab.~\ref{tab:crispdm} states what each phase meant concretely here, and where the project
departed from the standard description.

\begin{table}[H]
  \centering
  \small
  \begin{tabularx}{\linewidth}{lXX}
    \toprule
    \textbf{Phase} & \textbf{In this project} & \textbf{Adaptation} \\
    \midrule
    Business Understanding & The four difficulties of
      Section~\ref{subsec:problem}; the Business and Data Science Objectives &
      Extended with an explicit safety requirement: a confident wrong answer is worse
      than no answer \\
    \addlinespace[2pt]
    Data Understanding & Provenance, acquisition characteristics, class distribution, and
      the statistical unit of each dataset &
      Far heavier than anticipated. The statistical unit was undocumented in every imaging
      dataset and had to be recovered \\
    \addlinespace[2pt]
    Data Preparation & Cleaning, taxonomy correction, preprocessing, and case-level
      grouping &
      Grouping is treated as a modelling decision, not a preprocessing step, because it
      determines what every later metric means \\
    \addlinespace[2pt]
    Modeling & The four perception modules and the reasoning layer &
      Split across two levels: individual models, and the contract through which they
      cooperate \\
    \addlinespace[2pt]
    Evaluation & Case-level metrics against baselines, calibration, error analysis and
      explainability &
      A calibration step was inserted between modelling and evaluation: the deliverable is
      not a prediction but a prediction with a trustworthy confidence \\
    \addlinespace[2pt]
    Deployment & --- &
      Out of scope for the internship. Replaced by integration and a prototype
      (Chapter~\ref{sec:integration}); no clinical deployment was undertaken \\
    \bottomrule
  \end{tabularx}
  \caption{The six CRISP-DM phases as applied in this project.}
  \label{tab:crispdm}
\end{table}

\begin{figure}[H]
  \centering
  \includegraphics[width=0.72\textwidth]{fig_crispdm.png}
  \caption{The CRISP-DM cycle as followed here. The return path from modelling and
    evaluation to data understanding, drawn in red, ran repeatedly and was the dominant
    route through the cycle rather than an exception. Deployment was outside the scope of
    the internship.}
  \label{fig:crispdm}
\end{figure}

\paragraph{The loop back to the data was the main path, not an exception.} CRISP-DM allows
a return from Evaluation to earlier phases; in this project that return was where most of
the progress happened, and Fig.~\ref{fig:crispdm} draws it accordingly. Three times, a
model was trained and evaluated before the evaluation itself revealed that the data had
been misunderstood --- an unrecognised grouping structure in one dataset, a class whose
label did not match its images in another, an input feature that encoded the outcome in a
third. In each case the fix belonged to Data Understanding, not to Modeling, and the model
had to be discarded and retrained. A linear reading of the process would have concluded
after the first evaluation and reported a number that could not be defended.

\subsubsection{Development phases}

The work was organised into sequential phases, each concluding with a working, evaluated
component. The ordering was chosen so that every phase produced something the next phase
depended on, rather than developing components in parallel and integrating them at the
end. The phases are listed in Tab.~\ref{tab:phases}.

\begin{table}[H]
  \centering
  \small
  \begin{tabularx}{\linewidth}{clXl}
    \toprule
    \textbf{Phase} & \textbf{Objective} & \textbf{Outcome} & \textbf{Status} \\
    \midrule
    1 & Triage Agent & Three-tier urgency classifier (low / medium / high), calibrated
      & Complete \\
    \addlinespace[2pt]
    2 & Perception modules & Lung, cardiac and gallbladder models with calibrated
      confidence & Complete \\
    \addlinespace[2pt]
    3 & Unified Ultrasound Agent & Organ routing, shared output contract, encounter
      storage & Complete \\
    \addlinespace[2pt]
    4 & Clinical State Builder & Standardised representation of all agent outputs
      & Complete \\
    \addlinespace[2pt]
    5 & Clinical Reasoning Agent & Safety layer, escalation policy, retrieval and
      reasoning & In progress \\
    \bottomrule
  \end{tabularx}
  \caption{Development phases of the project.}
  \label{tab:phases}
\end{table}

\subsubsection{Methodological rules}

\paragraph{One methodological rule governed all phases.} No design decision was adopted on
the basis of expectation alone. Where a change was proposed --- a different architecture,
an augmentation strategy, an additional dataset --- it was implemented, measured against
the existing baseline, and retained only if the measurement supported it.

Several proposals that appeared well motivated were rejected on this basis, and those
rejections are reported in the chapters that follow rather than omitted. They are, in
several cases, more informative than the successes: a change that ought to have helped and
did not tells the reader something about the data that a favourable number does not. In a
clinical context, a result that cannot be defended is of no use regardless of how
favourable it appears.

\paragraph{Reporting convention.} Throughout Chapters~\ref{sec:perception}
and~\ref{sec:results}, three distinct statements are kept separate: what was
\textit{built}, what was \textit{measured}, and what may reasonably be \textit{concluded}.
A model architecture is a construction; a Dice coefficient on a held-out partition is a
measurement; a statement about how the model will behave on data from another hospital is
a conclusion, and requires more than the measurement to support it. Conflating the three
is the most common way a report overstates its results, and it is avoided here
deliberately.
